\uuid{KTrp}
\exo7id{6913}
\titre{exo7 6913}
\auteur{ruette}
\organisation{exo7}
\datecreate{2013-01-24}
\isIndication{false}
\isCorrection{false}
\chapitre{Probabilité discrète}
\sousChapitre{Variable aléatoire discrète}
\cpge{mpsi-mp2i.probabilites-univers-fini.d, pcsi.probabilites-variables-aleatoires.d, ptsi.probabilites-variables-aleatoires.d, bcpst1.variable-aleatoire-univers-fini.b, mp-mpi.variables-aleatoires-discretes.c, pc.probabilites-lois-usuelles.c, psi.probabilites-lois-usuelles.c, pt.probabilites-lois-usuelles.d, bcpst2.couples-variables-discretes.a}
\module{Probabilité et statistique}
\niveau{L2}
\difficulte{}

\contenu{
\texte{
Le tableau ci-dessous donne la loi d'un couple de variables aléatoires
$(X,Y)$, avec $X$ et $Y$ prenant chacune leurs valeurs dans $\{1,2,3\}$.
$$
\begin{array}{|c|c|c|c|}
\hline X\ \backslash\ Y & 1 & 2 & 3 \\ \hline
1&0&1/9&2/9\\ \hline
2&2/9&0&1/9\\ \hline
3&1/9&2/9&0\\
\hline\end{array}
$$
}
\begin{enumerate}
    \item \question{Calculer la probabilité que $X$ et $Y$ soit égales.}
    \item \question{Déterminer les lois de $X$ et de $Y$.}
\end{enumerate}
}
